\paperchapter{January 2002}{20 January 2002}{B}{50}{GSET 2002 (I)}
\keytable{\kcell{1}{A} & \kcell{2}{D} & \kcell{3}{D} & \kcell{4}{D} & \kcell{5}{C} & \kcell{6}{D} & \kcell{7}{C} & \kcell{8}{D} & \kcell{9}{C} & \kcell{10}{A} \\
\kcell{11}{C} & \kcell{12}{C} & \kcell{13}{B} & \kcell{14}{C} & \kcell{15}{D} & \kcell{16}{D} & \kcell{17}{B} & \kcell{18}{B} & \kcell{19}{D} & \kcell{20}{C} \\
\kcell{21}{D} & \kcell{22}{B} & \kcell{23}{C} & \kcell{24}{A} & \kcell{25}{A} & \kcell{26}{C} & \kcell{27}{A} & \kcell{28}{B} & \kcell{29}{B} & \kcell{30}{D} \\
\kcell{31}{C} & \kcell{32}{C} & \kcell{33}{A} & \kcell{34}{A} & \kcell{35}{D} & \kcell{36}{A} & \kcell{37}{A} & \kcell{38}{D} & \kcell{39}{B} & \kcell{40}{A} \\
\kcell{41}{D} & \kcell{42}{B} & \kcell{43}{B} & \kcell{44}{B} & \kcell{45}{D} & \kcell{46}{C} & \kcell{47}{C} & \kcell{48}{A} & \kcell{49}{C} & \kcell{50}{B}}
\begin{question}{1}
{Doctor and patient have same relationship as:\par}
\optstwo{Teacher and student}{Library and books}{Class and learning}{Husband and wife}
\answer{A}{Teacher and student}
\begin{explain}
A doctor serves and treats a patient; a \textbf{teacher} serves and teaches a \textbf{student}. The relationship is \emph{professional : the person served}.
\end{explain}
\end{question}

\begin{question}{2}
{Which one is different from the rest three ?\par}
\optsfour{Apple}{Orange}{Banana}{Fruit}
\answer{D}{Fruit}
\begin{explain}
Apple, orange and banana are \emph{kinds} of fruit; \textbf{fruit} is the general class, so it is different.
\end{explain}
\end{question}

\begin{question}{3}
{ACFJ: OUBJ:: SUXB: ?\par}
\optsfour{GNSA}{GLQZ}{GKPY}{GMTB}
\answer{D}{GMTB}
\begin{explain}
ACFJ: A(1), C(3), F(6), J(10), with gaps $+2, +3, +4$. OUBJ continues the same letters: O(15, $+5$), U(21, $+6$), B(28$\to$2, $+7$), J(10, $+8$). So the second group continues the first.
SUXB: S(19), U(21), X(24), B(28$\to$2). Continue: G(7, $+5$), M(13, $+6$), T(20, $+7$), B(28$\to$2, $+8$), giving \textbf{GMTB}.
\end{explain}
\end{question}

\begin{question}{4}
{BEAT: GIDV:: SOUP: ?\par}
\optsfour{XSYS}{YSXR}{YXSR}{XSXR}
\answer{D}{XSXR}
\begin{explain}
BEAT $\rightarrow$ GIDV: B$+5$, E$+4$, A$+3$, T$+2$.
SOUP $\rightarrow$ S$+5=$X, O$+4=$S, U$+3=$X, P$+2=$R, giving \textbf{XSXR}.
\end{explain}
\end{question}

\begin{question}{5}
{Which is the most effective method to control noise in a classroom?\par}
\optsone{Raising your voice above the students voice}{Saying 'don't talk'}{Remaining calm and just looking at students}{Continue teaching the class}
\answer{C}{Remaining calm and just looking at students}
\begin{explain}
Remaining calm and \textbf{silently looking at the students} usually restores order without confrontation. Shouting or scolding adds to the noise.
\end{explain}
\end{question}

\begin{question}{6}
{Choose one of the following statements as the most comprehensive way to define the job of a teacher:\par}
\optsone{Propagation of knowledge}{Conducting research}{Community service}{Propagation of knowledge, collection of knowledge and creation of knowledge}
\answer{D}{Propagation of knowledge, collection of knowledge and creation of knowledge}
\begin{explain}
A teacher's job comprises \textbf{propagation, collection (preservation) and creation of knowledge}, i.e.\ teaching, scholarship and research together.
\end{explain}
\end{question}

\begin{question}{7}
{A good academic researcher is:\par}
\optstwo{dependent on his teacher}{follower of research plan strictly}{open minded and radical}{highly knowledgable}
\answer{C}{open minded and radical}
\begin{explain}
A good researcher is \textbf{open-minded and radical}: willing to question established ideas and accept new evidence.
\end{explain}
\end{question}

\begin{question}{8}
{An effective teacher is one who:\par}
\optstwo{is very disciplinarian}{has mastery over his subject}{knows little but talk very well}{can motivate students to learn}
\answer{D}{can motivate students to learn}
\begin{explain}
An effective teacher \textbf{motivates students to learn}.
\end{explain}
\end{question}

\begin{question}{9}
{Symbiosis is beneficial to the :\par}
\optstwo{Host}{Parasite}{Host and parasite}{Neither host nor parasite}
\answer{C}{Host and parasite}
\begin{explain}
\textbf{Symbiosis} (mutualism) is an association in which \textbf{both partners benefit}, e.g.\ lichens (alga $+$ fungus). (The options loosely call them host and parasite.)
\end{explain}
\end{question}

\begin{question}{10}
{The disease 'AIDS' is chiefly due to :\par}
\optsone{infection in blood cells responsible for production of antibodies}{infection in cells of brain}{infection in cells of liver}{infection in cells of kidney}
\answer{A}{infection in blood cells responsible for production of antibodies}
\begin{explain}
HIV, the virus that causes AIDS, infects \textbf{helper T-lymphocytes (CD4 cells)}, which are essential for triggering antibody production and immune responses. This destroys the body's immunity.
\end{explain}
\end{question}

\begin{question}{11}
{Pasteurization of milk is done to :\par}
\optstwo{decrease the fat content}{increase the nutritive value}{destroy micro-organisms}{remove adulteration}
\answer{C}{destroy micro-organisms}
\begin{explain}
\textbf{Pasteurisation} (Louis Pasteur) heats milk (e.g.\ $72^\circ$C for 15 s, then rapid cooling) to \textbf{destroy harmful micro-organisms}.
\end{explain}
\end{question}

\begin{question}{12}
{Genetic engineering means :\par}
\optstwo{Mapping of DNA sequences}{Analysis of DNA structure}{Using recombinant DNA technique}{Understanding the functions of DNA}
\answer{C}{Using recombinant DNA technique}
\begin{explain}
Genetic engineering means deliberately modifying genes \textbf{using recombinant DNA (rDNA) technology}: cutting and joining DNA from different sources.
\end{explain}
\end{question}

\begin{context}
\textbf{Read the following passage and answer question nos.~13 to 16.}\par\smallskip
From the very beginning man has attempted what has seemed impossible. Man is different from the rest of the creation in this respect. He has an eternal thirst for adventure. This has led to countless new discoveries and inventions. Human curiosity is limitless. It is responsible for space flights and moon landings.

The desire to know what lies beyond the visible world takes many forms. The everest hero Tenzing and the hero of the seven Seas, Mihir Sen, were inspired by the same restless spirit. Astronaut Armstrong and his colleagues, who were the first humans to set foot on the soil of the moon, have proved beyond doubt that man shall not rest until he has conquered the entire universe.

Now the question arises: Is it enough to know and master nature? To answer it one should first answer a more basic question: which is more important --- knowing and understanding the world around you or knowing and understanding yourself? Everyone will agree that in the absence of self-knowledge even the most advanced knowledge of the universe will be not only useless but also dangerous.
\end{context}
\begin{question}{13}
{It is responsible for space flights and moon landings. Here 'It' means:\par}
\optstwo{Man's desire for adventure}{Man's curiosity}{Man's attempting the impossible}{Man and the rest of creation}
\answer{B}{Man's curiosity}
\begin{explain}
``Human curiosity is limitless. \textbf{It} is responsible for space flights and moon landings.'' \emph{It} refers to \textbf{man's curiosity}.
\end{explain}
\end{question}

\begin{question}{14}
{\ldots{}that man shall not rest until he has conquered the entire universe. Here 'rest' means:\par}
\optsfour{Sleep}{Sit and read}{Stop}{Lie down}
\answer{C}{Stop}
\begin{explain}
``Man shall not rest'' means man shall not \textbf{stop}.
\end{explain}
\end{question}

\begin{question}{15}
{He has an eternal thirst for adventure. Here 'eternal' means:\par}
\optstwo{Great}{Old}{Without beginning or end}{That which cannot be satisfied}
\answer{D}{That which cannot be satisfied}
\begin{explain}
Per the key, ``eternal thirst'' is a thirst \textbf{that cannot be satisfied} (a never-ending, insatiable thirst). (Literally, \emph{eternal} means without beginning or end.)
\end{explain}
\end{question}

\begin{question}{16}
{Man is different from the rest of the creation in this respect. Here 'Man' means:\par}
\optstwo{a particular man}{Tenzing and Mihir Sen}{Armstrong}{All mankind}
\answer{D}{All mankind}
\begin{explain}
``Man is different from the rest of the creation'' refers to \textbf{all mankind} (the human species).
\end{explain}
\end{question}

\begin{context}
\textbf{Read the following passage and answer question nos.~17 to 20.}\par\smallskip
The word medieval means `the middle age'. It is used in history, therefore, to refer to that period which lies between the ancient and the modern periods and is quite literally the middle period. How do we know when the ancient period ends and the medieval period begins? We have taken the eighth century A.D.\ as the beginning and the eighteenth century A.D.\ as the end of the medieval period. Why? Because as you will see when you read this book, there were many changes taking place in Indian Society in about the eighth century A.D.\ and during this time these changes influenced many aspects of Indian Society. They influenced the political and economic aspects, social laws, religion, language, art in short, almost everything. So we recognise that a new stage in Indian history had come about. We may say that this change took place around the eighth century.

We have to be just a little vague about the date because just as a person does not change suddenly in one year, similarly societies also take time to change. Nor does everybody become influenced by new ideas at the same time. Some of these changes in Indian history began earlier than the eighth century and in some parts of India their influence was felt a little later. But taking a general view of things, we can say that the new phase took root in the eighth century. In the same way with the break-up of the Mughal Empire and the coming of the British, the eighteenth century also saw many changes. So we refer to the century as the closing of the medieval period.
\end{context}
\begin{question}{17}
{We have to be just a little vague about the date.\ldots{} The word 'vague' here means :\par}
\optsfour{doubtful}{not precise}{misleading}{not clear}
\answer{B}{not precise}
\begin{explain}
``Vague about the date'' means \textbf{not precise}: the date cannot be fixed exactly, because change is gradual.
\end{explain}
\end{question}

\begin{question}{18}
{The eighth century has been taken as the beginning and the eighteenth century has been taken as the end of medieval period because :\par}
\optsone{No changes took place in Indian Society before eighth and after eighteenth century.}{A new stage began in Indian Society in the eighth and lasted up to eighteenth century.}{The period between eighth and eighteenth century is the most important period in Indian History.}{The Mughals came to India in the eighth and the British came to India in the eighteenth century.}
\answer{B}{A new stage began in Indian Society in the eighth and lasted up to eighteenth century.}
\begin{explain}
The author says that around the eighth century ``a new stage in Indian history had come about'', which ended with the changes of the eighteenth century. So \textbf{a new stage began in the eighth century and lasted up to the eighteenth}.
\end{explain}
\end{question}

\begin{question}{19}
{It is used in history.and is quite literally the middle period. The word 'literally' means:\par}
\optstwo{Approximately}{Strictly}{According to literature}{Accurately}
\answer{D}{Accurately}
\begin{explain}
``Quite literally the middle period'' means the word, taken in its exact sense, means ``middle period''. \emph{Literally} here means \textbf{accurately} (exactly).
\end{explain}
\end{question}

\begin{question}{20}
{The above extract has been taken from:\par}
\optstwo{a newspaper}{a history journal}{a book}{a magazine}
\answer{C}{a book}
\begin{explain}
``\ldots as you will see when you \textbf{read this book}'' shows the extract is from \textbf{a book}.
\end{explain}
\end{question}

\begin{question}{21}
{Human personality is a product of:\par}
\optstwo{Environment}{Heredity}{Society}{All the above three}
\answer{D}{All the above three}
\begin{explain}
Personality develops from the interaction of \textbf{heredity, environment and society}: all three.
\end{explain}
\end{question}

\begin{question}{22}
{Amnesia is known to mean :\par}
\optstwo{Sleep walking}{Loss of memory}{A symptom of stomach disease}{Alpha waves}
\answer{B}{Loss of memory}
\begin{explain}
\textbf{Amnesia} is \textbf{loss of memory}.
\end{explain}
\end{question}

\begin{question}{23}
{Long-term memory means:\par}
\optsone{Memory of long passages}{Exact location of a long road}{Remembering an event of long past}{Recalling a long material of immediate past}
\answer{C}{Remembering an event of long past}
\begin{explain}
Long-term memory stores information over long periods, e.g.\ \textbf{remembering an event of the long past}.
\end{explain}
\end{question}

\begin{question}{24}
{A man saves 5\% of his salary and pays 15\% of his salary towards taxes; rest of the salary is spent for family. If taxes paid by him are Rs. 30,000, find his salary.\par}
\optsfour{Rs. 2,00,000}{Rs. 3,00,000}{Rs. 1,50,000}{Rs. 1,00,000}
\answer{A}{Rs. 2,00,000}
\begin{explain}
Tax $=15\%$ of salary $=$ Rs.~30,000, so salary $=30{,}000\div0.15=$ \textbf{Rs.~2,00,000}.
\end{explain}
\end{question}

\begin{question}{25}
{If 2/5 of a number is equal to 2/3 of another number, what is the ratio between these two numbers ?\par}
\optsfour{5:3}{3:4}{15:4}{4:3}
\answer{A}{5:3}
\begin{explain}
$\tfrac25a=\tfrac23b\Rightarrow\dfrac ab=\dfrac{2/3}{2/5}=\dfrac53$. The ratio is \textbf{5 : 3}.
\end{explain}
\end{question}

\begin{question}{26}
{A 20'x14' wall is to be painted leaving out a window on it measuring 7'x4'. If the rate of painting is Rs. 10 per square foot, what will be the cost of painting ?\par}
\optsfour{2800}{2828}{2520}{3080}
\answer{C}{2520}
\begin{explain}
Area to paint $=20\times14-7\times4=280-28=252$ sq ft. Cost $=252\times10=$ \textbf{Rs.~2,520}.
\end{explain}
\end{question}

\begin{question}{27}
{Average of four numbers is 12. If two other numbers, whose average is 15 are added to the given collection, the resulting six numbers will have an average =\par}
\optsfour{13}{12}{15}{6}
\answer{A}{13}
\begin{explain}
Sum $=4\times12+2\times15=48+30=78$; average $=78\div6=\mathbf{13}$.
\end{explain}
\end{question}

\begin{question}{28}
{The last day of the year 2001 was a Monday. What will the first day of the year 2003 be?\par}
\optsfour{Monday}{Wednesday}{Tuesday}{Sunday}
\answer{B}{Wednesday}
\begin{explain}
31 December 2001 is Monday, so 1 January 2002 is Tuesday. 2002 has 365 days ($=52$ weeks $+1$ day), so 1 January 2003 is one day later: \textbf{Wednesday}.
\end{explain}
\end{question}

\begin{question}{29}
{I want to distribute Rs. 600 amongst Priya, Kanu and Geeta. If Priya receives 1/2 of the total amount and Kanu receives 1/3 of the total amount, what amount will Geeta receive ?\par}
\optsfour{200}{100}{400}{500}
\answer{B}{100}
\begin{explain}
Priya $=300$, Kanu $=200$, so Geeta $=600-500=$ \textbf{Rs.~100}.
\end{explain}
\end{question}

\begin{question}{30}
{A bedsheet costs Rs. 400. The shopowner offers a discount of 15\% on it. What is the net cost of the bedsheet ?\par}
\optsfour{415}{385}{460}{340}
\answer{D}{340}
\begin{explain}
Net cost $=400\times0.85=$ \textbf{Rs.~340}.
\end{explain}
\end{question}

\begin{question}{31}
{Which of the following is the correct meaning of 'Satyagraha' in Gandhian Thought ?\par}
\optstwo{Retaliation}{Aggression}{Discussion and Negotiation}{Military preparation}
\answer{C}{Discussion and Negotiation}
\begin{explain}
\textbf{Satyagraha} (``truth-force'') is Gandhi's method of non-violent resistance: persuading the opponent through truth, suffering, dialogue and negotiation rather than force. Among the options, \textbf{discussion and negotiation} is closest (per the key).
\end{explain}
\end{question}

\begin{question}{32}
{Attraction causes infatuation. Infatuation causes anger. Anger causes frustration. Therefore attraction causes frustration. Point out the correct statement from the following:\par}
\optstwo{Infatuation causes attraction.}{Frustration causes anger.}{Infatuation causes frustration.}{Anger causes Attraction.}
\answer{C}{Infatuation causes frustration.}
\begin{explain}
The chain is attraction $\rightarrow$ infatuation $\rightarrow$ anger $\rightarrow$ frustration. By the same chain, \textbf{infatuation causes frustration} (via anger). The other options reverse the direction of causation.
\end{explain}
\end{question}

\begin{question}{33}
{Gandhi and Einstein were contemporaries in Twentieth century. Gandhi was a social reformer. During Twentieth century there were many social reformers. Therefore Einstein was also a social reformer. Select from the following the statement which is correct:\par}
\optsone{Gandhi was a social reformer.}{All who lived in Twentieth century were social reformers.}{Einstein was a social reformer.}{Gandhi lived in Eighteenth century.}
\answer{A}{Gandhi was a social reformer.}
\begin{explain}
The argument is fallacious (Einstein does not become a reformer merely by living in the same century). The only correct statement among the options is the given premise, \textbf{``Gandhi was a social reformer''}.
\end{explain}
\end{question}

\begin{question}{34}
{M is the son of P. P and S are brothers. T is the mother of S. V is the daughter of S. Select the correct statement from the following:\par}
\optstwo{S is the uncle of M.}{T is the grand-mother of S.}{P is the sister of S.}{M and S are brothers.}
\answer{A}{S is the uncle of M.}
\begin{explain}
P and S are brothers and M is P's son, so \textbf{S is M's uncle}. (T is the mother of S, not the grandmother; P is a brother, not a sister; M and S are nephew and uncle.)
\end{explain}
\end{question}

\begin{question}{35}
{What is the wavelength-range of visible light radiations in electro-magnetic waves ?\par}
\optstwo{10\textasciicircum{}-14 meter to 10\textasciicircum{}-10 meter}{10\textasciicircum{}-5 meter to 10\textasciicircum{}-4 meter}{10\textasciicircum{}-6 meter to 10\textasciicircum{}-8 meter}{None of the above}
\answer{D}{None of the above}
\begin{explain}
Visible light spans about $4\times10^{-7}$ m to $7\times10^{-7}$ m (400--700 nm). None of the ranges given matches, so \textbf{none of the above}.
\end{explain}
\end{question}

\begin{question}{36}
{What is the state of mercury?\par}
\optsfour{Compound}{Metal}{Non-Metal}{Semi-conductor}
\answer{A}{Compound}
\keynote{Mercury is a \textbf{metal} (the only metal that is liquid at room temperature), which is option (B). The official key gives (A); answer (B) in your own understanding.}
\begin{explain}
Mercury (Hg) is an element, a \textbf{metal} that is liquid at room temperature. It is not a compound. The official key's letter (A) appears to be a misprint.
\end{explain}
\end{question}

\begin{question}{37}
{Which physical quantity has a unit kilo-watt-hour ?\par}
\optsfour{Power}{Energy}{Force}{Momentum}
\answer{A}{Power}
\keynote{The kilowatt-hour is a unit of \textbf{energy} (option B). The official key gives (A).}
\begin{explain}
$1\ \text{kWh}=1000\ \text{W}\times3600\ \text{s}=3.6\times10^6$ J. It is a unit of \textbf{energy} (the ``unit'' of electricity billing), not of power. The official key's (A) appears to be a misprint.
\end{explain}
\end{question}

\begin{question}{38}
{Which is the Unconventional source of energy ?\par}
\optsfour{Wood}{Coal}{Mineral Oil}{Tidal Energy}
\answer{D}{Tidal Energy}
\begin{explain}
\textbf{Tidal energy} is a non-conventional (renewable) source. Wood, coal and mineral oil are conventional.
\end{explain}
\end{question}

\begin{question}{39}
{What is the use of Naphtha ?\par}
\optstwo{Energy}{Chemical in a plant}{Solvent for petrochemicals}{Solvent for Biochemical}
\answer{B}{Chemical in a plant}
\begin{explain}
Naphtha is a light petroleum fraction used as a feedstock and solvent in the chemical and petrochemical industry. The official key gives \textbf{(B) chemical in a plant}.
\end{explain}
\end{question}

\begin{question}{40}
{What is the type of energy acquired by a body when lifted to certain height ?\par}
\optsfour{Kinetic}{Potential}{Heat}{Electrical}
\answer{A}{Kinetic}
\keynote{A body lifted to a height gains \textbf{potential energy} ($mgh$), option (B). The official key gives (A).}
\begin{explain}
Work done in lifting a body against gravity is stored as \textbf{gravitational potential energy} ($E_p=mgh$). The key's (A) Kinetic appears to be a misprint.
\end{explain}
\end{question}

\begin{question}{41}
{National Missile Defence Plan is undertaken by:\par}
\optstwo{Bill Clinton}{George Bush (Sr.)}{Ronald Reagan}{George Bush (Jr.)}
\answer{D}{George Bush (Jr.)}
\begin{explain}
The \textbf{National Missile Defense} programme was pushed forward by President \textbf{George W.~Bush (Jr.)}, who withdrew from the ABM Treaty in 2001--02.
\end{explain}
\end{question}

\begin{question}{42}
{Economic liberalization in India was initiated by:\par}
\optstwo{Rajiv Gandhi}{Narsimh Rao}{V. P. Singh}{Atal Bihari Vajpayee}
\answer{B}{Narsimh Rao}
\begin{explain}
Economic liberalisation (1991) was initiated by the \textbf{P.~V.~Narasimha Rao} government, with Dr.~Manmohan Singh as Finance Minister.
\end{explain}
\end{question}

\begin{question}{43}
{India exploded its first nuclear device in :\par}
\optsfour{1956}{1974}{1981}{1998}
\answer{B}{1974}
\begin{explain}
India's first nuclear test, \emph{Smiling Buddha}, was at Pokhran on 18 May \textbf{1974}.
\end{explain}
\end{question}

\begin{question}{44}
{Bank Nationalization and abolition of Privy Purses was initiated by:\par}
\optstwo{Jawaharlal Nehru}{Indira Gandhi}{Lal Bahadur Shastri}{Chandrashekhar}
\answer{B}{Indira Gandhi}
\begin{explain}
Bank nationalisation (1969) and the abolition of privy purses (26th Amendment, 1971) were initiated by \textbf{Indira Gandhi}.
\end{explain}
\end{question}

\begin{question}{45}
{The person who brought about integration of Princely states with the rest of India was:\par}
\optstwo{Morarji Desai}{Ram Manohar Lohia}{Vijaya Laxmi Pandit}{Sardar Patel}
\answer{D}{Sardar Patel}
\begin{explain}
\textbf{Sardar Vallabhbhai Patel}, as Home Minister, integrated over 560 princely states into the Indian Union.
\end{explain}
\end{question}

\begin{question}{46}
{The first Governor General of India was:\par}
\optstwo{Lord Mountbatten}{Dr. Rajendra Prasad}{C. Rajgopalachari}{Gen. Cariappa}
\answer{C}{C. Rajgopalachari}
\keynote{Lord Mountbatten was the first Governor-General of independent India (1947--48); C.~Rajagopalachari was the first and only \emph{Indian} Governor-General (1948--50). The official key gives (C).}
\begin{explain}
The key's intended answer is \textbf{C.~Rajagopalachari}, the first Indian to hold the office of Governor-General.
\end{explain}
\end{question}

\begin{question}{47}
{56 students are divided into three groups A, B and C. Group B consists of twice the number of students in group C. Group A has half the number of students than group C. How many students are there in group C ?\par}
\optsfour{8}{12}{16}{28}
\answer{C}{16}
\begin{explain}
Let group C $=x$. Then B $=2x$ and A $=x/2$, so $x+2x+\tfrac x2=56\Rightarrow3.5x=56\Rightarrow x=\mathbf{16}$.
\end{explain}
\end{question}

\begin{question}{48}
{You are walking straight, then you turned towards left. If the setting sun is now at your left, then in which direction were you walking in the beginning ?\par}
\optsfour{East}{West}{North}{South}
\answer{A}{East}
\begin{explain}
The setting sun is in the west. After turning left, the west is on your left, so you now face \textbf{north}. Before the left turn you were facing \textbf{east}.
\end{explain}
\end{question}

\begin{question}{49}
{While going to office, you find that a running car has just run over a child. What will you do?\par}
\optsone{Hire a taxi and try to follow the car}{Inform the nearest police station about the accident}{Render first-aid to the injured child}{Search the pockets of the injured to find out his home address}
\answer{C}{Render first-aid to the injured child}
\begin{explain}
The first duty is to save life: \textbf{render first-aid} to the injured child (then seek medical help and inform the police).
\end{explain}
\end{question}

\begin{question}{50}
{Index of mental ability is known as:\par}
\optsfour{EQ}{IQ}{MQ}{SQ}
\answer{B}{IQ}
\begin{explain}
Mental ability is expressed as the \textbf{Intelligence Quotient (IQ)} $=\dfrac{\text{Mental age}}{\text{Chronological age}}\times100$.
\end{explain}
\end{question}
